\section{Supplementary material}\label{app:supp}

This appendix records three deferred studies of the rough rock joint of Section~\ref{sec:cv7}. Each
study isolates one question that the main text raises but does not develop. The first (S.1) asks whether
the transition-map detector survives being wired into the contact loop of an unsteady solve, rather than
only queried on a fixed point cloud. The second (S.2) asks whether the dilatant interface obeys the first
law over a full load cycle, with the energy ledger closed term by term. The third (S.3) asks whether an
independent explicit solver reproduces the same frictional response. We report each study with its
measured numbers and with the limit it exposes. The level set (a neural signed-distance function, or SDF)
runs throughout as the comparison baseline, on the same geometry and the same query set as the chart.

\subsection{Transition-map detection inside the contact loop}\label{app:supp:p3}

The main text reports the chart and the level set as static geometry estimators, each queried on a fixed
grid. Here we instead let the transition-map detector supply the gap and the normal that drive a
quasi-static direct-shear solve, so the detector is evaluated inside the Newton loop. The chart reads
contact as the boundary-to-boundary map $\tmap=\map{B}^{-1}\!\circ\map{A}$, which sends a point on one
face to the mating point on the other: for a query point on face $A$, the height chart returns the mating
surface height on face $B$ directly, and the gap then follows as the signed normal separation between the
two. The level set returns the same gap as the value of the ambient neural SDF $\phi$, with normal
$\nbf=\grad\phi/\norm{\grad\phi}$.

% code: cv7_transition_map_contact.py::_analytic
Throughout, $\gn$ denotes the signed normal gap, with the convention $\gn<0$ in penetration. The detector
declares a query point in contact when its gap is negative, so the active set is the index set of those
points,
%
\begin{equation}\label{eq:supp_activeset}
\mathcal{A}=\{\,p : \gn(\xbf_p)<0\,\},
\end{equation}
%
where $p$ indexes the query points and $\xbf_p$ is the position of point $p$. To score the active set, we
compare it against the analytic active set $\mathcal{A}^{\star}$, which we obtain from the closed-form
mating surface. The score is the symmetric agreement
%
\begin{equation}\label{eq:supp_agreement}
\eta=1-\frac{\abs{\mathcal{A}\,\triangle\,\mathcal{A}^{\star}}}{N},
\end{equation}
%
in which $\triangle$ is the symmetric difference (the points the two sets disagree on) and $N$ is the
number of query points; $\eta=1$ means the detector and the closed form agree on every point. We measure
gap accuracy separately, as the relative root-mean-square error of the predicted gap $\gn$ against the
analytic gap $\gn^{\star}$,
%
\begin{equation}\label{eq:supp_gaprmse}
e_g=\frac{\norm{\gn-\gn^{\star}}_{2}}{\norm{\gn^{\star}}_{2}}.
\end{equation}
%
The gap error matters because the gap, not the active-set membership, sets the contact pressure.

% code: cv7_transition_map_contact.py::_chart cv7_transition_map_contact.py::_sdf
On $N=10^{5}$ query points sampled over the joint footprint, the chart detector reproduced the analytic
active set to $\eta=98.9\%$, against $95.8\%$ for the level set. On the gap, the chart was sharper by a
wider margin: $e_g=4.2\%$ for the chart against $44.7\%$ for the level set, a factor of $10.5$ on the
quantity that sets the contact pressure. Wired into the FEM shear, the chart-driven solve tracked the
analytic-surface solve closely, reaching a peak apparent friction $\muapp=0.566$ against the analytic
$0.513$, with the shear traction agreeing to within $9.5\%$ over the loading history. Table~\ref{tab:supp_p3}
collects these results.

% code: cv7_transition_map_contact.py::ang
One quantity runs the other way. At the asperity tips the level set carries the sharper
normal, $0.8^{\circ}$ of deviation against $10.3^{\circ}$ for the chart, because the random-Fourier height
field has a noisier gradient at the tips than the smooth ambient field does. The chart advantage is on the
gap and the active set, the two quantities that decide where contact occurs and how hard, not on the
pointwise normal direction at the sharpest features. We state the tip-normal result so the comparison is
not read as uniform: the level set is a fair baseline, and on the tip-normal direction it is the more
accurate estimator on this surface.

% code: cv7_transition_map_contact.py::main
The chart detector costs $1.47\times$ the level set per query. A chart query is a single forward
evaluation of $\map{A}$ followed by the chart read of $\map{B}^{-1}$, with no level-set root find or ray
march, because the mating height is the chart output. The factor below two is the overhead of evaluating
two Fourier-feature charts in place of one ambient SDF, and it stays constant as the number of query
points grows. The shear traction and the emergent normal dilation that this detector produces over the
loading history appear in Figure~\ref{fig:supp_traction}.

\begin{table}
\centering
\caption{Transition-map detector wired into the FEM contact loop (Phase~3), scored against the analytic
mating surface on $N=10^{5}$ query points. Gap $\gn<0$ denotes penetration. The chart wins on active-set
agreement and gap error; the level set has the sharper asperity-tip normal on this random-Fourier surface.}
\label{tab:supp_p3}
\begin{tabular}{lccc}
\toprule
quantity & chart & level set (SDF) & note \\
\midrule
active-set agreement $\eta$ & $98.9\%$ & $95.8\%$ & Eq.~\eqref{eq:supp_agreement} \\
gap error $e_g$            & $4.2\%$  & $44.7\%$ & Eq.~\eqref{eq:supp_gaprmse}; $10.5\times$ sharper \\
tip-normal deviation       & $10.3^{\circ}$ & $0.8^{\circ}$ & level set sharper \\
peak $\muapp$ (FEM shear)  & $0.566$  & $0.513^{\dagger}$ & $^{\dagger}$analytic-surface reference \\
shear traction $\tau$      & \multicolumn{2}{c}{within $9.5\%$ of analytic} & chart-driven vs analytic \\
per-query cost             & $1.47\times$ & $1\times$ & constant in $N$ \\
\bottomrule
\end{tabular}
\end{table}

\begin{figure}
\centering
\includegraphics[width=\linewidth]{cv7_traction_history_pub}
\caption{Phase~3 transition-map detection inside the direct-shear loop. Global shear traction $\tau$ and
emergent normal dilation $\delta_n$ as functions of the imposed tangential displacement, with the chart
detector supplying the gap and normal to the FEM. The traces are mesh-converged across the refinement
sequence of Section~\ref{sec:cv7}; the chart-driven response tracks the analytic-surface reference, and
Table~\ref{tab:supp_p3} reports the quantitative agreement.}
\label{fig:supp_traction}
\end{figure}

\subsection{Cyclic loading and the energy ledger}\label{app:supp:p4}

The cyclic study drives the joint through a triangle-wave tangential schedule under constant normal
displacement. It accounts for every term of the first law. The friction model uses a stateful
return-map update, so stick and slip alternate and the shear-traction response forms hysteresis loops.
The Plesha asperity-degradation parameter decays across cycles, from $0.3405$ to $0.3377$ over four
cycles, and the contact Newton residual averages $7\times10^{-5}$. The peak apparent friction is
$\abs{\muapp}\approx0.33$.

% code: rock_joint_decoder_cyclic.py::run_cyclic
The first law states that the external work fed into the joint must reappear as stored energy plus
dissipation. With $\gn<0$ in penetration and slip taken as plastic, the ledger reads
%
\begin{equation}\label{eq:supp_ledger}
W_{\mathrm{ext}}
= \Delta U_{\mathrm{el}}
+ W_{\mathrm{pen}}
+ W_{\mathrm{stick}}
+ W_{\mathrm{fric}},
\end{equation}
%
where each term has a physical meaning. The left side, $W_{\mathrm{ext}}=\sum \fbf_{\mathrm{contact}}\cdot
\mathrm{d}\ubf_{\mathrm{platen}}$, is the machine work the contact tractions do on the moving platen. On
the right, $\Delta U_{\mathrm{el}}=\tfrac12\ubf^{\!\top}\!K\ubf$ is the recoverable strain energy of the
blocks, $W_{\mathrm{pen}}=\tfrac12\sum\varepsilon_n A\,\gn^{2}$ is the recoverable normal penalty energy,
$W_{\mathrm{stick}}$ is the recoverable tangential spring energy of the stuck nodes, and
$W_{\mathrm{fric}}=\sum\abs{\tbf_T}\,\abs{\mathrm{d}\bm{s}_p}\ge0$ is the frictional dissipation
accumulated over the plastic slip increments. The three recoverable terms make up the moving part of the
balance over a cycle: each one returns to its start value once the load closes,
$\Delta U_{\mathrm{el}}\to0$, $W_{\mathrm{pen}}\to0$, $W_{\mathrm{stick}}\to0$. The per-cycle closure
therefore reduces to a single ratio that should equal one,
%
\begin{equation}\label{eq:supp_closure}
\Pi_{\mathrm{cyc}}
= \frac{W_{\mathrm{ext}}}
{\,W_{\mathrm{fric}}+\Delta U_{\mathrm{el}}+W_{\mathrm{pen}}+W_{\mathrm{stick}}\,}
\;\longrightarrow\; 1 .
\end{equation}
%
A value of $\Pi_{\mathrm{cyc}}=1$ means the work in matches the energy out, term for term, over the cycle.

% code: rock_joint_decoder_cyclic.py::energy_report
Measured per cycle, $\Pi_{\mathrm{cyc}}=1.09$ on the first cycle and $1.01$ on the second, the latter
inside the acceptance band $[0.98,1.02]$. The first-cycle excess is the recoverable terms not yet having
returned to zero after a single traversal; by the second cycle they do return, and the ledger closes to
one percent. Figure~\ref{fig:supp_cyclic} shows the hysteresis loops and the per-cycle ledger.

% code: rock_joint_decoder_cyclic.py::energy_report
The \emph{cumulative} ratio over the whole history is a different number, $0.17$, far from one. The
shear-only constant-normal machine work $W_{\mathrm{ext}}$ does not count the dissipation spent seating
the asperities under the initial normal load. The run begins with
$W_{\mathrm{fric}}[0]=0.91$ already accumulated while $W_{\mathrm{ext}}[0]=0$, because the normal-load
seating happens before the tangential schedule starts. That offset enters the denominator of the
cumulative ratio but never the numerator, so it depresses the ratio permanently. The offset is a
bookkeeping artifact, the unmeasured seating dissipation, not an energy violation. The per-cycle balance
of Eq.~\eqref{eq:supp_closure} differences the accumulated work across one full cycle and so cancels the
constant offset; that per-cycle balance is the correct first-law statement, and it closes.

\begin{figure}
\centering
\includegraphics[width=\linewidth]{rock_joint_cyclic_energy_pub}
\caption{Phase~4 cyclic shear under constant normal displacement. Left: stick/slip hysteresis loops of the
shear traction against tangential displacement, with Plesha degradation decaying across cycles. Right: the
energy ledger of Eq.~\eqref{eq:supp_ledger}; the per-cycle closure $\Pi_{\mathrm{cyc}}$ of
Eq.~\eqref{eq:supp_closure} is $1.09$ then $1.01$. The cumulative ratio $0.17$ is the uncounted initial
asperity-seating dissipation, a bookkeeping offset rather than an energy violation.}
\label{fig:supp_cyclic}
\end{figure}

\subsection{Explicit ChartMPM cross-check}\label{app:supp:p5}

% code: rock_joint_mpm_xcheck.py::RoughFloorSDF
The last study checks the frictional response with an independent solver. We discretize a deformable block
with the explicit finite-deformation ChartMPM (linear B-spline transfers, Neo-Hookean response) and shear
it against the \emph{same} band-limited rough surface that the chart-FEM uses, sliding the surface beneath
the block at a constant rate to impose direct-shear kinematics. The contact runs through the existing
per-particle MPM channel: the gap comes from \code{evaluate\_gap}, the normal force from the penalty model,
and the tangential force from regularized Coulomb friction, all scattered as physical-space per-particle
forces through the same projection-to-grid channel as gravity. The solver sees the floor as the analytic
SDF $\phi(\xbf)=z-h(x-u_x,y)$, where $h$ is the rough height field and $u_x$ is the imposed slide.

% code: rock_joint_mpm_xcheck.py::main
The explicit solve reproduced the Coulomb friction floor. A flat-bottomed block sliding on the rough floor
reached an emergent apparent friction $\muapp\approx0.34$ against the prescribed base $\mu=0.4$, within
about $14\%$, and converged stably under dynamic relaxation. This is an independent dynamic confirmation
of the frictional channel that the quasi-static FEM uses.

% code: rock_joint_mpm_xcheck.py::RoughFloorSDF
The explicit solve does not, however, confirm the dilatant interlock. The chart-FEM produces emergent
dilation and an apparent-friction enhancement above the base value because two rough faces ride over one
another. A \emph{mated} rough-on-rough block in the explicit solver, with the full surface in conformal
contact, is unstable in explicit penalty dynamics and diverges without an implicit or mortar contact
treatment. The cross-check therefore validates friction but not dilatancy. We state this limit rather
than hide it: the explicit MPM is a consistent check of the Coulomb response, and the emergent dilatancy
of Section~\ref{sec:cv7} rests instead on the quasi-static implicit chart-FEM, whose consistent tangent
drives the contact Newton residual to $\sim10^{-9}$. An implicit or mortar contact treatment for the
mated MPM case is the documented path to a full dynamic cross-check of dilatancy.
